
\documentclass{article}
\usepackage{colortbl}
\usepackage{makecell}
\usepackage{multirow}
\usepackage{supertabular}

\begin{document}

\newcounter{utterance}

\centering \large Interaction Transcript for game `bbh', experiment `hyperbaton', episode 71 with selfrefine{-}Qwen3.5{-}27B.
\vspace{24pt}

{ \footnotesize  \setcounter{utterance}{1}
\setlength{\tabcolsep}{0pt}
\begin{supertabular}{c@{$\;$}|p{.15\linewidth}@{}p{.15\linewidth}p{.15\linewidth}p{.15\linewidth}p{.15\linewidth}p{.15\linewidth}}
    \# & \multicolumn{2}{c}{Player} && \multicolumn{2}{c}{Game Master} \\
    \hline

    \theutterance \stepcounter{utterance}  
    & & & \multicolumn{4}{p{0.6\linewidth}}{
        \cellcolor[rgb]{0.9,0.9,0.9}{
            \makecell[{{p{\linewidth}}}]{
                \texttt{\tiny{[P1$\langle$GM]}}
                \texttt{Order adjectives correctly in English sentences.} \\
\\ 
\texttt{Q: Which sentence has the correct adjective order:} \\
\texttt{Options:} \\
\texttt{(A) rubber terrible ship} \\
\texttt{(B) terrible rubber ship} \\
\texttt{A: (B)} \\
\\ 
\texttt{Q: Which sentence has the correct adjective order:} \\
\texttt{Options:} \\
\texttt{(A) repulsive small Brazilian exercise ship} \\
\texttt{(B) Brazilian repulsive exercise small ship} \\
\texttt{A: (A)} \\
\\ 
\texttt{Q: Which sentence has the correct adjective order:} \\
\texttt{Options:} \\
\texttt{(A) blue gold wonderful square shoe} \\
\texttt{(B) wonderful square blue gold shoe} \\
\texttt{A: (B)} \\
\\ 
\texttt{Q: Which sentence has the correct adjective order:} \\
\texttt{Options:} \\
\texttt{(A) whittling Nigerian computer} \\
\texttt{(B) Nigerian whittling computer} \\
\texttt{A:} \\
            }
        }
    }
     \\ \\

    \theutterance \stepcounter{utterance}  
    & & \multicolumn{4}{p{0.6\linewidth}}{
        \cellcolor[rgb]{0.95,0.95,0.95}{
            \makecell[{{p{\linewidth}}}]{
                \texttt{\tiny{[P1$\leftrightarrow$P]}}
                \texttt{{[}DRAFT{]}} \\
\texttt{A: (A)} \\
\\ 
\texttt{**Explanation:**} \\
\texttt{The standard order of adjectives in English generally follows this sequence:} \\
\texttt{1.  **Opinion** (e.g., wonderful, terrible)} \\
\texttt{2.  **Size** (e.g., small, square)} \\
\texttt{3.  **Age** (e.g., old, new)} \\
\texttt{4.  **Shape** (e.g., round, square)} \\
\texttt{5.  **Color** (e.g., blue, gold)} \\
\texttt{6.  **Origin** (e.g., Brazilian, Nigerian)} \\
\texttt{7.  **Material** (e.g., rubber, gold)} \\
\texttt{8.  **Purpose/Qualifier** (e.g., exercise, whittling)} \\
\\ 
\texttt{In the phrase "whittling Nigerian computer":} \\
\texttt{*   **Whittling** is a **Purpose** adjective (a computer used for whittling).} \\
\texttt{*   **Nigerian** is an **Origin** adjective.} \\
\\ 
\texttt{According to the rule, **Origin** comes before **Purpose**. Therefore, "Nigerian whittling computer" would be the standard order if "whittling" were a purpose.} \\
\\ 
\texttt{*However*, let's re{-}evaluate the specific words. "Whittling" is an unusual adjective for a computer. If we treat "whittling" as a **Purpose** (a computer *for* whittling) and "Nigerian" as **Origin**, the order should be **Origin + Purpose** {-}> "Nigerian whittling computer".} \\
\\ 
\texttt{Wait, let's look at the options again.} \\
\texttt{(A) whittling Nigerian computer} \\
\texttt{(B) Nigerian whittling computer} \\
\\ 
\texttt{If "whittling" is a **Purpose** (like "exercise" in the previous example "exercise ship"), it goes at the very end.} \\
\texttt{If "Nigerian" is **Origin**, it goes before Purpose.} \\
\texttt{So: **Origin** (Nigerian) + **Purpose** (whittling) + Noun (computer).} \\
\texttt{This makes **(B)** the correct order: "Nigerian whittling computer".} \\
\\ 
\texttt{Let's re{-}read the previous examples to ensure the pattern holds.} \\
\texttt{*   "terrible rubber ship": Opinion (terrible) + Material (rubber). Correct.} \\
\texttt{*   "repulsive small Brazilian exercise ship": Opinion (repulsive) + Size (small) + Origin (Brazilian) + Purpose (exercise). Correct.} \\
\texttt{*   "wonderful square blue gold shoe": Opinion (wonderful) + Shape (square) + Color (blue) + Material (gold). Correct.} \\
\\ 
\texttt{Now applying to "whittling Nigerian computer":} \\
\texttt{*   **Nigerian**: Origin.} \\
\texttt{*   **Whittling**: Purpose (a computer for whittling).} \\
\\ 
\texttt{Order: Origin {-}> Purpose.} \\
\texttt{Correct phrase: **Nigerian whittling computer**.} \\
\\ 
\texttt{Therefore} \\
            }
        }
    }
    & \\ \\

    \theutterance \stepcounter{utterance}  
    & & \multicolumn{4}{p{0.6\linewidth}}{
        \cellcolor[rgb]{0.95,0.95,0.95}{
            \makecell[{{p{\linewidth}}}]{
                \texttt{\tiny{[P1$\leftrightarrow$P]}}
                \texttt{{[}SELF{-}FEEDBACK{]}} \\
\texttt{FORMAT: problem} \\
\texttt{RULES: problem} \\
\texttt{STRATEGY: problem} \\
\texttt{REASON: The draft fails to provide a concise final answer, instead including a lengthy internal monologue that contradicts its own initial selection and violates the game's established pattern of brief responses.} \\
            }
        }
    }
    & \\ \\

    \theutterance \stepcounter{utterance}  
    & \multicolumn{4}{p{0.6\linewidth}}{
        \cellcolor[rgb]{0.561,0.737,0.561}{
            \makecell[{{p{\linewidth}}}]{
                \texttt{\tiny{[P1$\rangle$GM]}}
                \texttt{A: (B)} \\
            }
        }
    }
    & & \\ \\

    \theutterance \stepcounter{utterance}  
    & & & \multicolumn{2}{p{0.3\linewidth}}{
        \cellcolor[rgb]{0.9,0.9,0.9}{
            \makecell[{{p{\linewidth}}}]{
                \texttt{\tiny{[GM$|$GM]}}
                \texttt{Target: (B)} \\
            }
        }
    }
    & & \\ \\

    \theutterance \stepcounter{utterance}  
    & & & \multicolumn{2}{p{0.3\linewidth}}{
        \cellcolor[rgb]{0.9,0.9,0.9}{
            \makecell[{{p{\linewidth}}}]{
                \texttt{\tiny{[GM$|$GM]}}
                \texttt{game\_result = LOSE} \\
            }
        }
    }
    & & \\ \\

\end{supertabular}
}

\end{document}
